% =====================================================================================
% estilo.tex: formatação do TCC segundo as Diretrizes para Elaboração de TCC de
% Graduação do IFTO (2. ed., 2020) e a ABNT NBR 14724, sobre a classe abnTeX2.
% Compila com XeLaTeX ou LuaLaTeX (recomendado) e também com pdfLaTeX.
% =====================================================================================

% ---------- Fonte: Arial 12 pt (TeX Gyre Heros, de métricas equivalentes às da Arial) ----------
\usepackage{iftex}
\ifPDFTeX
  \usepackage[T1]{fontenc}
  \usepackage[utf8]{inputenc}
  \usepackage{tgheros}
  \renewcommand*{\familydefault}{\sfdefault}
\else
  \usepackage{fontspec}
  \setmainfont{texgyreheros}[Extension=.otf, UprightFont=*-regular, BoldFont=*-bold,
    ItalicFont=*-italic, BoldItalicFont=*-bolditalic]
  \setsansfont{texgyreheros}[Extension=.otf, UprightFont=*-regular, BoldFont=*-bold,
    ItalicFont=*-italic, BoldItalicFont=*-bolditalic]
\fi

% ---------- Pacotes ----------
\usepackage{indentfirst}          % recuo também no primeiro parágrafo de cada seção
\usepackage{graphicx}
\usepackage{array}
\usepackage{longtable}            % quadros e tabelas que podem passar de uma página
\usepackage{enumitem}
\usepackage{amsmath,amssymb}
\usepackage{xurl}                 % endereços longos quebram em qualquer ponto
\usepackage{pdfpages}             % para incluir a ficha catalográfica em PDF
\urlstyle{same}                   % endereço na mesma fonte do texto, como pede a ABNT

% ---------- Página: A4, margens 3 cm (superior e esquerda) e 2 cm (inferior e direita) ----------
\setlrmarginsandblock{3cm}{2cm}{*}
\setulmarginsandblock{3cm}{2cm}{*}
\checkandfixthelayout

% ---------- Número da página: só o número, no canto superior direito, em 10 pt ----------
\makeoddhead{abntheadings}{}{}{\footnotesize\thepage}
\makeevenhead{abntheadings}{\footnotesize\thepage}{}{}
\makeheadrule{abntheadings}{\textwidth}{0pt}

% ---------- Parágrafos: recuo de 1,25 cm, espaço 1,5, sem espaço extra entre parágrafos ----------
\setlength{\parindent}{1.25cm}
\setlength{\parskip}{0pt}
\setlength{\emergencystretch}{3em}   % evita linhas estouradas com nomes de código longos

% ---------- Títulos: todos em 12 pt; primário em CAIXA ALTA e negrito, secundário em negrito,
%            terciário sem destaque, quaternário em itálico (como na versão em Word) ----------
\def\ABNTEXchapterfont{\bfseries}
\def\ABNTEXchapterfontsize{\normalsize}
\def\ABNTEXsectionfont{\bfseries}
\def\ABNTEXsectionfontsize{\normalsize}
\def\ABNTEXsubsectionfont{\normalfont}
\def\ABNTEXsubsectionfontsize{\normalsize}
\def\ABNTEXsubsubsectionfont{\itshape}
\def\ABNTEXsubsubsectionfontsize{\normalsize}

% Sumário com o mesmo destaque dos títulos no texto (NBR 6027)
\renewcommand{\cftchapterfont}{\bfseries}
\renewcommand{\cftchapterpagefont}{\bfseries}
\renewcommand{\cftsectionfont}{\bfseries}
\renewcommand{\cftsectionpagefont}{\bfseries}
\renewcommand{\cftsubsectionfont}{\normalfont}
\renewcommand{\cftsubsectionpagefont}{\normalfont}

% Título sem número, centralizado e no sumário (REFERÊNCIAS, APÊNDICE A...)
\newcommand{\capitulosemnumero}[1]{%
  \begingroup
    \renewcommand{\printchaptertitle}[1]{\centering\ABNTEXchapterfont\ABNTEXchapterfontsize ##1}%
    \chapter*{#1}%
  \endgroup
  \addcontentsline{toc}{chapter}{\protect\numberline{}#1}}

% ---------- Nomes das listas ----------
\addto\captionsbrazil{%
  \renewcommand{\listfigurename}{Lista de figuras}%
  \renewcommand{\listtablename}{Lista de tabelas}%
  \renewcommand{\figurename}{Figura}%
  \renewcommand{\tablename}{Tabela}}

% ---------- Gráficos e quadros: contadores e listas próprios (receita do abnTeX2) ----------
\newcommand{\graficoname}{Gráfico}
\newcommand{\listofgraficosname}{Lista de gráficos}
\newfloat[chapter]{grafico}{lgf}{\graficoname}      % extensão lgf: "log" é o log do LaTeX
\newlistof{listofgraficos}{lgf}{\listofgraficosname}
\newlistentry{grafico}{lgf}{0}
\counterwithout{grafico}{chapter}

\newcommand{\quadroname}{Quadro}
\newcommand{\listofquadrosname}{Lista de quadros}
\newfloat[chapter]{quadro}{loq}{\quadroname}
\newlistof{listofquadros}{loq}{\listofquadrosname}
\newlistentry{quadro}{loq}{0}
\counterwithout{quadro}{chapter}

\counterwithout{figure}{chapter}
\counterwithout{table}{chapter}
\counterwithout{equation}{chapter}   % equação (1), e não (4.1)

% Entradas das listas: "Figura 1 — Título ...... página"
\renewcommand{\cftfigurename}{\figurename\space}
\renewcommand{\cfttablename}{\tablename\space}
\renewcommand{\cftgraficoname}{\graficoname\space}
\renewcommand{\cftquadroname}{\quadroname\space}
\renewcommand*{\cftfigureaftersnum}{\hfill\textemdash\hfill}
\renewcommand*{\cfttableaftersnum}{\hfill\textemdash\hfill}
\renewcommand*{\cftgraficoaftersnum}{\hfill\textemdash\hfill}
\renewcommand*{\cftquadroaftersnum}{\hfill\textemdash\hfill}
\setlength{\cftfigurenumwidth}{6.2em}
\setlength{\cfttablenumwidth}{6.6em}
\setlength{\cftgraficonumwidth}{6.8em}
\setlength{\cftquadronumwidth}{7.4em}
\setlength{\cftfigureindent}{0pt}
\setlength{\cfttableindent}{0pt}
\setlength{\cftgraficoindent}{0pt}
\setlength{\cftquadroindent}{0pt}

% ---------- Legendas: acima da ilustração, centralizadas, em 10 pt; fonte abaixo, em 10 pt ----------
\makeatletter
\newcommand{\legenda}[2]{% #1 = contador (figure, grafico, quadro, table); #2 = título
  \refstepcounter{#1}%
  \addcontentsline{\csname ext@#1\endcsname}{#1}%
    {\protect\numberline{\csname the#1\endcsname}{\ignorespaces #2}}%
  {\centering\SingleSpacing\footnotesize
   \csname #1name\endcsname~\csname the#1\endcsname~\textemdash~#2\par}}
\makeatother
\providecommand{\fonte}{}
\renewcommand{\fonte}[1]{\par\vspace{4pt}{\centering\SingleSpacing\footnotesize Fonte: #1\par}}

% Figura ou gráfico: legenda, imagem e fonte ficam juntas, no ponto do texto ou no pé da página
% (posição [hbp]): nunca antes do trecho que as cita e sem deixar meia página em branco.
\newcommand{\legendafig}[3]{\centering\legenda{#1}{#2}\label{#3}\vspace{4pt}}
\renewcommand{\topfraction}{0.9}
\renewcommand{\bottomfraction}{0.8}
\renewcommand{\textfraction}{0.07}
\renewcommand{\floatpagefraction}{0.75}
\setcounter{topnumber}{2}
\setcounter{bottomnumber}{2}
\setcounter{totalnumber}{4}

% Quadros e tabelas (longtable): a legenda é a primeira linha e a fonte, a última,
% para que nunca se separem da tabela; o cabeçalho se repete nas páginas seguintes.
\newlength{\largtab}
\newcolumntype{L}[1]{>{\raggedright\arraybackslash\hspace{0pt}}p{#1\largtab}}
\newcolumntype{C}[1]{>{\centering\arraybackslash\hspace{0pt}}p{#1\largtab}}
% O longtable soma 1 ao contador de tabelas a cada uso; o -1 desfaz isso, e só \legenda numera.
\newenvironment{tabelaIFTO}
  {\par\addtocounter{table}{-1}\SingleSpacing\footnotesize\setlength{\extrarowheight}{2pt}%
   \setlength{\LTpre}{0.5\baselineskip}\setlength{\LTpost}{0.5\baselineskip}}
  {\par}
\newcommand{\legendatab}[4]{% #1 contador, #2 título, #3 rótulo, #4 colunas
  \multicolumn{#4}{@{}c@{}}{\parbox{\linewidth}{\legenda{#1}{#2}\label{#3}\vspace{4pt}}}}
% Quadro ou tabela curta: um bloco só, que vai inteiro para a página seguinte se não couber
\newenvironment{tabelacurta}
  {\par\addvspace{0.5\baselineskip}\noindent\begin{minipage}{\linewidth}\centering
   \SingleSpacing\footnotesize\setlength{\extrarowheight}{2pt}}
  {\end{minipage}\par\addvspace{0.5\baselineskip}}
\newcommand{\fontetab}[2]{% #1 colunas, #2 texto da fonte
  \multicolumn{#1}{@{}c@{}}{\parbox{\linewidth}{\vspace{4pt}\centering\footnotesize Fonte: #2}}}

% ---------- Alíneas (NBR 6024), passos numerados e lista de conferência ----------
\newlist{alineasIFTO}{enumerate}{1}
\setlist[alineasIFTO]{label=\alph*), labelindent=\parindent, leftmargin=*, align=left,
  topsep=0pt, partopsep=0pt, itemsep=0pt, parsep=0pt}
\newlist{passos}{enumerate}{1}
\setlist[passos]{label=\arabic*., labelindent=\parindent, leftmargin=*, align=left,
  topsep=0pt, partopsep=0pt, itemsep=0pt, parsep=0pt}
\newlist{conferir}{itemize}{1}
\setlist[conferir]{label=$\square$, labelindent=\parindent, leftmargin=*, align=left,
  topsep=0pt, partopsep=0pt, itemsep=0pt, parsep=0pt}

% ---------- Referências: espaço simples, alinhadas à esquerda, separadas por uma linha em branco ----------
\newenvironment{referenciasIFTO}
  {\par\SingleSpacing\raggedright\setlength{\parindent}{0pt}\setlength{\parskip}{\baselineskip}}
  {\par}

% ---------- Campos ainda não preenchidos ----------
% (o \kern impede que a fonte junte << e >> em « e »)
\newcommand{\pendente}[1]{\textless\kern0pt\textless#1\textgreater\kern0pt\textgreater}
